% =============================================================================
% 02-computer-system-organization.tex — Week 2:
% Computer-System Organization and Architecture
%
% Course: Operating Systems (PCC CS-502)
% Anchor reading: Silberschatz, Galvin & Gagne, Operating System Concepts,
%                 8th ed. — Ch. 1 §1.2 (computer-system organization,
%                 interrupts, storage, I/O, DMA, pp. 6-18), §1.3
%                 (computer-system architecture, pp. 15-20), §1.7-1.8
%                 (memory and storage management, pp. 24-26); page-verified
%                 via .../SilberschatzOSC/index.md.
%                 Stallings, Operating Systems: Internals and Design
%                 Principles, 6th ed. — Ch. 1 §1.4 (interrupts, p. 15),
%                 §1.5 (memory hierarchy, p. 26), §1.6 (cache, p. 29),
%                 §1.7 (I/O techniques and DMA, p. 33); page-verified via
%                 .../StallingsOS/index.md.
%                 Bovet & Cesati, Understanding the Linux Kernel, 3rd ed. —
%                 Ch. 4 (interrupts and exceptions, pp. 131-188), optional
%                 concrete-kernel depth; page-verified via .../LinuxKernel/index.md.
%
% Compile with: /compile-latex CS502/02-computer-system-organization
% =============================================================================

\documentclass[aspectratio=169]{beamer}
\input{header}   % resolves via TEXINPUTS=../../Preambles;$TEXINPUTS — compile
                  % from cwd Slides/CS502/ (TeX writes output next to cwd)
\usepackage{array}
\coursecode{CS502}

\title{Computer-System Organization and Architecture}
\subtitle{Week 2 --- Interrupts, Storage Structure, I/O Structure, and DMA}
\author[Dr. Auqib Hamid Lone]{%
  \textbf{Dr. Auqib Hamid Lone}\\
  \small Assistant Professor in Computer Science \& Engineering
}
\institute[GCET Ganderbal]{%
  \small\textbf{PCC CS-502: Operating Systems}\\[0.3em]
  Department of Computer Science \& Engineering\\
  Government College of Engineering \& Technology (GCET), Ganderbal
}
\date[\today]{\small\today}

\begin{document}

\frame{\titlepage}

% =============================================================================
% ACT 0 --- THE MACHINE THE OS MUST MANAGE
% =============================================================================
\section{The Machine the OS Must Manage}

\begin{frame}{Welcome: Five Events Before One Character Appears}
  \begin{block}{The big question}
    You press \texttt{Enter} to run \texttt{cat bigfile.txt}. Before a single
    byte reaches the screen, the keyboard controller, the CPU, the memory, the
    disk controller, and the bus must all cooperate. Which of them can signal
    the CPU --- and how?
  \end{block}
  \begin{itemize}
    \item The OS does not run ``on its own''. It is \emph{event-driven}: it
      reacts to hardware events it did not initiate.
    \item Last week: what the OS \emph{is}. This week: the machine it stands on,
      and the two mechanisms --- \key{interrupts} and \key{DMA} --- that make it
      possible to stand there.
  \end{itemize}
  \vspace{0.15cm}
  \begin{alertblock}{Running thread for today}
    We follow the command \texttt{cat bigfile.txt} from the Enter key to the
    disk platter and back: an interrupt starts it, a storage hierarchy feeds it,
    and DMA moves the bytes.
  \end{alertblock}
\end{frame}

\begin{frame}{Roadmap: Five Questions About the Machine}
  \begin{enumerate}
    \item \key{What is the machine made of?} --- CPU, memory, controllers, bus.
    \item \key{How does hardware get the CPU's attention?} --- interrupts,
      traps, and the interrupt vector.
    \item \key{Where do data live?} --- main memory, disks, SSDs, and the
      memory hierarchy.
    \item \key{How do bytes move?} --- polling, interrupt-driven I/O, and DMA.
    \item \key{How does the machine scale?} --- multiprocessors, multicore, and
      clusters.
  \end{enumerate}
  \vspace{0.2cm}
  \muted{Every answer is an OS-visible fact: each mechanism in this lecture
  reappears later as a scheduling, memory, or I/O policy.}
\end{frame}

\transitionslide{A computer is not a CPU that does everything. It is a set of independent devices that run concurrently and communicate over a shared bus.}

% =============================================================================
% ACT 1 --- COMPUTER-SYSTEM ORGANIZATION
% =============================================================================
\section{Computer-System Organization}

\begin{frame}{Anatomy of a Computer System}
  \centering
  \begin{tikzpicture}[every node/.style={font=\small}]
    \node[draw=primary-blue, fill=light-bg, minimum width=3.0cm,
          minimum height=1.0cm] (cpu) at (-4.2,1.55) {CPU};
    \node[draw=primary-blue, fill=light-bg, minimum width=3.0cm,
          minimum height=1.0cm] (mem) at (0,1.55) {Main memory};
    \node[draw=primary-blue, fill=light-bg, minimum width=3.0cm,
          minimum height=1.0cm] (disk) at (4.2,1.55) {Disk controller};
    \node[draw=primary-gold, fill=white, minimum width=12.6cm,
          minimum height=0.55cm] (bus) at (0,0) {System bus (address, data, control)};
    \node[draw=primary-blue, fill=light-bg, minimum width=3.0cm,
          minimum height=1.0cm] (usb) at (-4.2,-1.55) {USB controller};
    \node[draw=primary-blue, fill=light-bg, minimum width=3.0cm,
          minimum height=1.0cm] (gpu) at (0,-1.55) {Graphics adapter};
    \node[draw=primary-blue, fill=light-bg, minimum width=3.0cm,
          minimum height=1.0cm] (net) at (4.2,-1.55) {Network controller};
    \draw[observed-edge] (cpu) -- (bus);
    \draw[observed-edge] (mem) -- (bus);
    \draw[observed-edge] (disk) -- (bus);
    \draw[observed-edge] (usb) -- (bus);
    \draw[observed-edge] (gpu) -- (bus);
    \draw[observed-edge] (net) -- (bus);
  \end{tikzpicture}
  \vspace{0.15cm}
  \begin{itemize}
    \item Each device has a \key{controller} --- hardware that speaks the bus
      protocol and operates the device.
    \item The CPU and controllers run \key{concurrently}; the bus is the shared
      resource they contend for.
  \end{itemize}
\end{frame}

\begin{frame}{Bootstrap: How the Machine Starts}
  \begin{enumerate}
    \item \key{Firmware} (BIOS/UEFI) runs from ROM/flash when power is applied.
      It tests the hardware (POST) and finds a boot device.
    \item The \key{bootstrap program} loads the OS kernel from disk into memory
      and jumps to its entry point.
    \item The kernel initializes its data structures, starts the first process
      (\texttt{init} / \texttt{systemd}), and from then on is \key{event-driven}.
  \end{enumerate}
  \vspace{0.2cm}
  \begin{alertblock}{Why the OS cannot assume it is running}
    Until the kernel is loaded, no OS exists. After it is loaded, the kernel
    only runs when a \emph{trap}, an \emph{interrupt}, or an explicit call gives
    it the CPU. The next section is about that handoff.
  \end{alertblock}
  \vspace{0.1cm}
  \muted{Bootstrap sequence: Silberschatz et al., Ch.~1 (pp.~6--8); Stallings,
  Ch.~1 \cite{StallingsOS_internals_design_principles}.}
\end{frame}

\begin{frame}{The CPU's View: Fetch, Decode, Execute --- and Wait}
  \begin{itemize}
    \item The CPU repeats one cycle: \key{fetch} an instruction at the program
      counter, \key{decode} it, \key{execute} it, update the PC.
    \item Registers hold the working set: general-purpose registers, the
      \key{program counter (PC)}, the stack pointer, and the status word.
    \item I/O is orders of magnitude slower than the CPU. A 3~GHz core executes
      roughly 3 instructions per nanosecond; a disk access takes milliseconds.
  \end{itemize}
  \vspace{0.2cm}
  \begin{block}{The utilization problem}
    If the CPU waited for every device, it would idle almost all the time.
    The OS needs a way to \emph{start} an I/O operation and be told later that it
    finished --- without watching it.
  \end{block}
  \vspace{0.1cm}
  \muted{Instruction cycle and registers: Stallings, Ch.~1 \S1.2--1.3
  (pp.~7--15) \cite{StallingsOS_internals_design_principles}.}
\end{frame}

\begin{frame}{Controllers and Drivers: Hardware and Its Software Twin}
  \begin{block}{Device controller (hardware)}
    A chip or board with registers the CPU can read and write: \key{status}
    (busy, ready, error), \key{control} (start, reset, mode), and
    \key{data-in}/\key{data-out} buffers. It operates the physical device.
  \end{block}
  \begin{block}{Device driver (software)}
    Kernel code that knows the controller's register layout and protocol. The
    OS calls the driver; the driver programs the controller; the controller
    moves the data.
  \end{block}
  \vspace{0.15cm}
  \begin{itemize}
    \item One driver per controller family; the rest of the kernel sees a
      uniform interface (read, write, ioctl).
    \item The kernel keeps a \key{device-status table} listing each device, its
      type, address, and state.
  \end{itemize}
  \vspace{0.1cm}
  \muted{Controllers, drivers, and the device-status table: Silberschatz et al.,
  Ch.~1 (pp.~6--18).}
\end{frame}

\begin{frame}{Worked Thread: The Enter Key, Step by Step}
  \begin{enumerate}
    \item You press \texttt{Enter}. The keyboard controller scans the key matrix
      and places a scan code in its data register.
    \item The controller asserts an \key{interrupt request (IRQ)} line on the bus.
    \item The CPU finishes the current instruction, saves its state, and jumps
      to the keyboard \key{interrupt service routine (ISR)}.
    \item The ISR reads the scan code, maps it to a character, and queues it for
      the terminal process.
    \item The scheduler eventually runs the shell, which reads the character and
      starts \texttt{cat bigfile.txt} --- which now needs the disk.
  \end{enumerate}
  \vspace{0.15cm}
  \begin{alertblock}{The pattern to remember}
    Device $\to$ controller $\to$ IRQ $\to$ CPU saves state $\to$ ISR $\to$
    kernel decides what happens next. Every device in the machine uses a variant
    of this path.
  \end{alertblock}
\end{frame}

\begin{frame}{Socratic Check 1}
  \begin{block}{Think, then answer}
    Why can the OS not simply poll every device in a loop, checking whether each
    one is ready? Name two distinct costs of that design.
  \end{block}
  \vspace{0.2cm}
  \begin{itemize}
    \item \key{CPU cost}: polling consumes cycles that could run user code ---
      and most polls find nothing ready.
    \item \key{Latency cost}: a slow or sporadic device (a keyboard, a network
      packet) may wait an entire polling period before being noticed.
    \item Interrupts invert the arrangement: the device tells the CPU when it
      needs attention, instead of the CPU asking.
  \end{itemize}
\end{frame}

\transitionslide{The CPU is fast and the devices are slow and independent. The only workable contract is: finish your instruction, then let the device interrupt you.}

% =============================================================================
% ACT 2 --- INTERRUPTS
% =============================================================================
\section{Interrupts and Traps}

\begin{frame}{Interrupts vs.\ Traps: Two Kinds of Transfer}
  \begin{block}{Interrupt (asynchronous)}
    Raised by \emph{hardware} at an unpredictable time: a key press, a packet
    arrival, a disk completing a transfer. It is unrelated to the instruction
    the CPU is currently executing.
  \end{block}
  \begin{block}{Trap or exception (synchronous)}
    Raised by the     \emph{CPU} because of the instruction just executed: a system
    call, a divide by zero, or an illegal opcode. It is caused by the program
    itself.
  \end{block}
  \vspace{0.15cm}
  \begin{itemize}
    \item Both transfer control to a kernel handler and both save the CPU state.
    \item The distinction matters: an interrupt can arrive \emph{between} any two
      instructions; a trap is tied to a specific instruction.
  \end{itemize}
\end{frame}

\begin{frame}{The Interrupt Path}
  \centering
  \begin{tikzpicture}[every node/.style={font=\small}]
    \node[draw=primary-blue, fill=light-bg, minimum width=2.5cm,
          minimum height=0.9cm] (dev) at (-4.8,1.0) {Device};
    \node[draw=primary-blue, fill=light-bg, minimum width=2.5cm,
          minimum height=0.9cm] (pic) at (-1.6,1.0) {Interrupt controller};
    \node[draw=primary-blue, fill=light-bg, minimum width=2.5cm,
          minimum height=0.9cm] (cpu) at (1.6,1.0) {CPU};
    \node[draw=primary-blue, fill=light-bg, minimum width=2.5cm,
          minimum height=0.9cm] (isr) at (4.8,1.0) {ISR};
    \node[draw=primary-gold, fill=white, minimum width=11.2cm,
          minimum height=0.75cm] (vec) at (0,-0.6)
      {Interrupt vector table: IRQ number $\to$ handler address};
    \draw[observed-edge] (dev) -- (pic);
    \draw[observed-edge] (pic) -- (cpu);
    \draw[observed-edge] (cpu) -- (isr);
    \draw[observed-edge] (cpu.south) -- (cpu.south |- vec.north);
  \end{tikzpicture}
  \vspace{0.2cm}
  \begin{enumerate}
    \item Device asserts IRQ; the interrupt controller prioritizes and forwards it.
    \item CPU finishes the current instruction (not the whole program), saves
      the PC and status word, and switches to kernel mode.
    \item CPU indexes the \key{interrupt vector table} by IRQ number and jumps
      to the handler.
    \item The ISR services the device; on return, the saved state is restored.
  \end{enumerate}
\end{frame}

\begin{frame}{The Interrupt Vector Table, Masking, and Priority}
  \begin{itemize}
    \item The \key{interrupt vector table} is an array of handler addresses,
      indexed by interrupt number. It lets the CPU dispatch without searching.
    \item \key{Maskable} interrupts can be temporarily disabled by the kernel
      (when updating shared structures). \key{Non-maskable} interrupts
      (power failure, fatal hardware error) cannot.
    \item \key{Priority levels} decide which interrupt wins when several arrive
      at once; a low-priority handler can itself be interrupted by a
      high-priority one.
  \end{itemize}
  \vspace{0.2cm}
  \begin{alertblock}{The design tension}
    Handlers must be short, or they block other interrupts and add latency. So
    real kernels split work: a minimal \key{top half} acknowledges the device and
    a deferred \key{bottom half} does the rest (Linux softirqs, tasklets, work
    queues).
  \end{alertblock}
  \vspace{0.1cm}
  \muted{Linux top-half/bottom-half handling: Bovet \& Cesati, Ch.~4
  (pp.~131--188) \cite{BovetCesati_understanding_linux_kernel}.}
\end{frame}

\begin{frame}{Worked Example: What Does Interrupt Handling Cost?}
  \begin{block}{Setup}
    Saving the context and entering the handler costs $2\ \mu\text{s}$; the
    handler itself runs for $8\ \mu\text{s}$. Each interrupt therefore costs
    $10\ \mu\text{s}$ of CPU time.
  \end{block}
  \begin{itemize}
    \item At \key{1{,}000 interrupts/s}: $1{,}000 \times 10\ \mu\text{s} =
      10\ \text{ms/s} = \good{1\%}$ CPU overhead.
    \item At \key{10{,}000 interrupts/s}: $10{,}000 \times 10\ \mu\text{s} =
      100\ \text{ms/s} = \bad{10\%}$ CPU overhead.
    \item At \key{100{,}000 interrupts/s}: the machine spends
      $\bad{100\%}$ of its time handling interrupts and does no useful work.
  \end{itemize}
  \vspace{0.2cm}
  \begin{alertblock}{The lesson}
    Interrupts are cheap per event but not free. When a device produces a
    \emph{stream} of events, per-event interrupts stop scaling --- which is
    exactly the problem DMA solves for block transfers.
  \end{alertblock}
\end{frame}

\begin{frame}{Socratic Check 2}
  \begin{block}{Think, then answer}
    An interrupt handler runs with interrupts temporarily disabled. The handler
    is long. What can go wrong, and what does this imply about handler design?
  \end{block}
  \vspace{0.2cm}
  \begin{itemize}
    \item Other devices cannot deliver their interrupts, so their data can be
      overwritten or lost --- a \key{missed interrupt}.
    \item The system's \key{interrupt latency} grows; a real-time device may miss
      its deadline.
    \item Hence the top-half/bottom-half split: acknowledge quickly, defer the
      heavy work to a context where interrupts are enabled.
  \end{itemize}
\end{frame}

\transitionslide{Interrupts tell the CPU that data has arrived. But where does the data live, and why can a program believe it has more memory than the machine owns?}

% =============================================================================
% ACT 3 --- STORAGE STRUCTURE AND THE MEMORY HIERARCHY
% =============================================================================
\section{Storage Structure and the Memory Hierarchy}

\begin{frame}{Main Memory and Secondary Storage}
  \begin{block}{Main memory (RAM)}
    A large array of bytes, each with an address, implemented in DRAM.
    \good{Fast} (tens of nanoseconds) and directly addressable by the CPU;
    \bad{volatile} --- its contents vanish when power is lost.
  \end{block}
  \begin{block}{Secondary storage}
    \key{Hard disks} (magnetic platters) and \key{solid-state drives}
    (NAND flash). \good{Non-volatile} and far larger than RAM, but
    \bad{orders of magnitude slower} --- microseconds for an SSD, milliseconds
    for a spinning disk.
  \end{block}
  \vspace{0.15cm}
  \begin{itemize}
    \item Programs live in memory only while they run; files persist on
      secondary storage.
    \item The OS moves data between the two in units of \key{blocks} (disk) or
      \key{cache lines} (memory).
  \end{itemize}
\end{frame}

\begin{frame}{Hard-Disk Geometry: Where Access Time Comes From}
  \centering
  \begin{tikzpicture}[every node/.style={font=\small}]
    \node[draw=primary-blue, fill=light-bg, minimum width=3.4cm,
          minimum height=1.0cm] (seek) at (-4.4,0.8) {Seek};
    \node[draw=primary-blue, fill=light-bg, minimum width=3.4cm,
          minimum height=1.0cm] (rot) at (0,0.8) {Rotational latency};
    \node[draw=primary-blue, fill=light-bg, minimum width=3.4cm,
          minimum height=1.0cm] (xfer) at (4.4,0.8) {Transfer};
    \node[draw=primary-gold, fill=white, minimum width=12.4cm,
          minimum height=0.8cm] (tot) at (0,-0.8)
      {Total access time $=$ seek $+$ rotational latency $+$ transfer};
    \draw[observed-edge] (seek) -- (tot.north -| seek.south);
    \draw[observed-edge] (rot) -- (tot.north -| rot.south);
    \draw[observed-edge] (xfer) -- (tot.north -| xfer.south);
  \end{tikzpicture}
  \vspace{0.15cm}
  \begin{itemize}
    \item \key{Platters} rotate; each surface has concentric \key{tracks} split
      into \key{sectors}; the same track across platters is a \key{cylinder}.
    \item \key{Seek}: move the arm to the track (mechanical, milliseconds).
    \item \key{Rotational latency}: wait for the sector to pass under the head
      (milliseconds).
    \item \key{Transfer}: stream the requested sectors (microseconds per block).
  \end{itemize}
\end{frame}

\begin{frame}{Worked Example: A 4 KB Disk Read at 7200 RPM}
  \begin{block}{Parameters}
    Seek time $= 8\ \text{ms}$; rotation speed $= 7{,}200\ \text{RPM}$;
    transfer rate $= 150\ \text{MB/s}$; block size $= 4\ \text{KB}$.
  \end{block}
  \begin{itemize}
    \item One rotation: $60/7{,}200 = \bad{8.33\ \text{ms}}$; average rotational
      latency $= 8.33/2 = \bad{4.17\ \text{ms}}$.
    \item Transfer time: $4\ \text{KB} / 150\ \text{MB/s} \approx
      \good{27\ \mu\text{s}}$.
    \item Total: $8 + 4.17 + 0.027 \approx \bad{12.2\ \text{ms}}$, of which
      \emph{99.8\%} is mechanical waiting.
  \end{itemize}
  \vspace{0.15cm}
  \begin{alertblock}{The consequence for the OS}
    Disk performance is dominated by \key{where the head is}, not by how fast
    bytes stream. This is why scheduling the \emph{order} of requests (later,
    disk scheduling) matters so much.
  \end{alertblock}
\end{frame}

\begin{frame}{The Memory Hierarchy}
  \centering
  \begin{tikzpicture}[every node/.style={font=\small}]
    \node[draw=primary-blue, fill=light-bg, minimum width=5.2cm,
          minimum height=0.65cm] (r) at (0,3.0) {Registers};
    \node[draw=primary-blue, fill=light-bg, minimum width=6.4cm,
          minimum height=0.65cm] (l1) at (0,2.25) {Cache (L1/L2/L3)};
    \node[draw=primary-blue, fill=light-bg, minimum width=7.6cm,
          minimum height=0.65cm] (m) at (0,1.5) {Main memory (DRAM)};
    \node[draw=primary-blue, fill=light-bg, minimum width=8.8cm,
          minimum height=0.65cm] (ssd) at (0,0.75) {Solid-state drive};
    \node[draw=primary-blue, fill=light-bg, minimum width=10.0cm,
          minimum height=0.65cm] (hdd) at (0,0) {Magnetic disk};
    \node[draw=primary-blue, fill=light-bg, minimum width=11.2cm,
          minimum height=0.65cm] (tape) at (0,-0.75) {Tape / optical archive};
  \end{tikzpicture}
  \vspace{0.15cm}
  \begin{itemize}
    \item Moving down: capacity grows, cost per byte falls, access time grows by
      orders of magnitude.
    \item Moving up: data is \key{cached} --- a copy is kept in the faster level
      on the bet that it will be used again.
  \end{itemize}
\end{frame}

\begin{frame}{Locality, Caching, and Why the Hierarchy Works}
  \begin{block}{The locality principle}
    Programs do not access memory uniformly. \key{Temporal locality}: an item
    used once is likely to be used again soon. \key{Spatial locality}: an item
    near one just used is likely to be used soon.
  \end{block}
  \begin{itemize}
    \item A \key{cache} is a small, fast store holding copies of recently used
      data. A \key{hit} is served from the cache; a \key{miss} goes to the next
      level and installs a copy.
    \item The OS manages caches at every level: CPU caches (hardware), the
      \key{buffer cache} for disk blocks, and memory itself as a cache for
      disk.
  \end{itemize}
  \vspace{0.1cm}
  \muted{Caching and the memory hierarchy: Stallings, Ch.~1 \S1.5--1.6
  (pp.~26--32) \cite{StallingsOS_internals_design_principles}; Silberschatz
  et al., Ch.~1 (pp.~9--14).}
\end{frame}

\begin{frame}{Socratic Check 3}
  \begin{block}{Think, then answer}
    Main memory is volatile, yet a file you saved last week still exists after
    a power cut. Where is the file, and what role did RAM play in saving it?
  \end{block}
  \vspace{0.2cm}
  \begin{itemize}
    \item The file lives on \key{secondary storage} (disk or SSD), which is
      non-volatile.
    \item RAM held the program's buffer and the kernel's \key{cache} while the
      write was in progress; the OS copied the dirty blocks to storage.
    \item \muted{This is why ``save'' is an I/O operation, not a memory
      operation --- and why a crash before the write-back can lose data.}
  \end{itemize}
\end{frame}

\transitionslide{Data must come from a device into memory. The question is who moves it: the CPU, one byte at a time, or a controller acting on its own.}

% =============================================================================
% ACT 4 --- I/O STRUCTURE: POLLING, INTERRUPTS, AND DMA
% =============================================================================
\section{I/O Structure: Polling, Interrupts, and DMA}

\begin{frame}{Three Ways to Move Data}
  \centering
  \small
  \setlength{\tabcolsep}{4pt}
  \begin{tabular}{@{}>{\raggedright\arraybackslash}p{2.2cm}
      >{\raggedright\arraybackslash}p{3.7cm}
      >{\raggedright\arraybackslash}p{3.7cm}
      >{\raggedright\arraybackslash}p{3.2cm}@{}}
    \toprule
    & \textbf{Polling} & \textbf{Interrupt-driven} & \textbf{DMA} \\
    \midrule
    Data mover & CPU & CPU & Controller \\
    CPU cost & Wasted while waiting & One handler per event & Setup $+$ one completion \\
    Best for & Fast, dedicated devices & Slow, sporadic devices & Block transfers \\
    Weakness & Busy-wait; latency & Overhead per event & Setup cost; bus contention \\
    \bottomrule
  \end{tabular}
  \vspace{0.35cm}
  \begin{itemize}
    \item The three techniques are a progression: the CPU does \emph{more}
      waiting, then \emph{less} work, then \emph{no} per-byte work.
    \item \texttt{cat bigfile.txt} uses a \key{block device}: DMA is the natural
      choice; a keyboard uses \key{interrupt-driven} I/O.
  \end{itemize}
\end{frame}

\begin{frame}{Polling: Simple but Wasteful}
  \begin{block}{Programmed I/O (polling)}
    The CPU repeatedly reads the controller's \key{status register} until the
    device is ready, then moves the data word by word.
  \end{block}
  \begin{itemize}
    \item \good{Simple}: no interrupt machinery, no concurrency.
    \item \bad{Wasteful}: every poll is a bus transaction and a wasted
      instruction slot; a slow device wastes almost all of them.
    \item \good{Still right} for very fast devices (a GPU register) or when the
      wait is predictably short --- polling a device that is always ready beats
      paying interrupt overhead.
  \end{itemize}
  \vspace{0.15cm}
  \muted{Programmed I/O: Silberschatz et al., Ch.~1 \S1.2.3 (pp.~12--18);
  Stallings, Ch.~1 \S1.7 (p.~33) \cite{StallingsOS_internals_design_principles}.}
\end{frame}

\begin{frame}{Interrupt-Driven I/O: The CPU Is Notified}
  \begin{enumerate}
    \item The driver starts the device by writing its control register.
    \item The CPU returns to other work; the scheduler runs another process.
    \item When the device finishes, it raises an interrupt.
    \item The ISR reads the data (or a status), acknowledges the device, and
      wakes the waiting process.
  \end{enumerate}
  \vspace{0.2cm}
  \begin{alertblock}{The cost model}
    CPU cost scales with the \emph{number of events}, not the number of bytes.
    Small blocks therefore mean many interrupts; large blocks mean few. This is
    the lever DMA pulls.
  \end{alertblock}
\end{frame}

\begin{frame}{Direct Memory Access (DMA)}
  \centering
  \begin{tikzpicture}[every node/.style={font=\small}]
    \node[draw=primary-blue, fill=light-bg, minimum width=2.6cm,
          minimum height=0.95cm] (cpu) at (-4.6,1.4) {CPU};
    \node[draw=primary-blue, fill=light-bg, minimum width=2.6cm,
          minimum height=0.95cm] (mem) at (0,1.4) {Main memory};
    \node[draw=primary-blue, fill=light-bg, minimum width=2.6cm,
          minimum height=0.95cm] (dma) at (0,-0.2) {DMA controller};
    \node[draw=primary-blue, fill=light-bg, minimum width=2.6cm,
          minimum height=0.95cm] (dc) at (4.6,1.4) {Disk controller};
    \draw[<->, thick, draw=jet] (cpu) -- (mem);
    \draw[<->, thick, draw=jet] (mem) -- (dc);
    \draw[<->, thick, draw=jet] (dma) -- (mem);
    \draw[<->, thick, draw=jet] (dma) -- (dc);
    \draw[observed-edge, dashed] (cpu) -- node[below, font=\scriptsize] {setup} (dma);
  \end{tikzpicture}
  \vspace{0.2cm}
  \begin{enumerate}
    \item The CPU programs the DMA controller with the memory address, the byte
      count, and the device.
    \item The controller transfers the whole block directly to or from memory,
      stealing bus cycles but no CPU instructions.
    \item The controller raises \emph{one} interrupt when the block is complete.
  \end{enumerate}
  \vspace{0.1cm}
  \muted{DMA: Silberschatz et al., Ch.~1 \S1.2.3; Stallings, Ch.~1 \S1.7
  (p.~33) \cite{StallingsOS_internals_design_principles}.}
\end{frame}

\begin{frame}{Worked Example: Moving 1 MB Three Ways}
  \begin{block}{Parameters}
    Disk transfer rate $= 200\ \text{MB/s}$; interrupt cost $= 10\ \mu\text{s}$;
    DMA setup $+$ completion interrupt $\approx 20\ \mu\text{s}$; transfer time
    for 1~MB $= 5\ \text{ms}$.
  \end{block}
  \begin{itemize}
    \item \key{Interrupt-driven, 4 KB blocks}: $256$ interrupts $\times
      10\ \mu\text{s} = \bad{2.56\ \text{ms}}$ --- about 51\% of the transfer.
    \item \key{Interrupt-driven, 64 KB blocks}: $16$ interrupts $\times
      10\ \mu\text{s} = \good{160\ \mu\text{s}}$ --- about 3.2\%.
    \item \key{DMA}: one setup plus one completion interrupt $\approx
      \good{20\ \mu\text{s}}$ --- about 0.4\%, while the CPU runs other work.
  \end{itemize}
  \vspace{0.15cm}
  \begin{alertblock}{The trade-off, stated precisely}
    Larger blocks amortize interrupt overhead but increase latency and buffer
    memory. DMA removes per-byte CPU work but adds setup cost --- so it wins on
    large transfers, not on a single byte.
  \end{alertblock}
\end{frame}

\begin{frame}{Device Types and the Device-Status Table}
  \begin{itemize}
    \item \key{Block devices}: addressable in fixed-size blocks, seekable ---
      disks, SSDs, USB drives. They support file systems and DMA.
    \item \key{Character devices}: a stream of bytes, not seekable ---
      keyboards, serial ports, mice. They are naturally interrupt-driven.
    \item The kernel maintains a \key{device-status table}: for each device, its
      type, address, and current state (idle, busy, error).
    \item When an interrupt arrives, the kernel looks up the device in the table,
      updates its state, and wakes any process waiting on it.
  \end{itemize}
  \vspace{0.2cm}
  \muted{Device types and the device-status table: Silberschatz et al., Ch.~1
  (pp.~12--18).}
\end{frame}

\begin{frame}{Socratic Check 4}
  \begin{block}{Think, then answer}
    DMA removes the CPU from the data path. Does that mean a DMA transfer uses
    no CPU time at all? Name the hidden costs.
  \end{block}
  \vspace{0.2cm}
  \begin{itemize}
    \item \key{Setup}: the CPU programs address, count, and direction for every
      transfer.
    \item \key{Bus contention}: the DMA controller steals memory cycles, slowing
      the CPU's own memory accesses (cycle stealing).
    \item \key{Cache coherence}: the CPU may hold stale copies of memory the DMA
      wrote; the OS and hardware must invalidate or flush.
    \item Conclusion: DMA is a \emph{trade} of per-byte CPU work for per-block
      setup and bus pressure --- not a free lunch.
  \end{itemize}
\end{frame}

\transitionslide{One CPU, one memory, one disk is the simple case. Real machines replicate processors and whole systems --- and the OS must keep them coherent.}

% =============================================================================
% ACT 5 --- COMPUTER-SYSTEM ARCHITECTURE
% =============================================================================
\section{Computer-System Architecture}

\begin{frame}{From One Processor to Many}
  \begin{block}{Single-processor systems}
    One general-purpose CPU. Most personal computers and phones have several
    cores per chip, but the OS still sees a set of processors to schedule.
  \end{block}
  \begin{itemize}
    \item \key{Multiprocessor} systems: two or more processors share the bus,
      clock, and memory.
      \begin{itemize}
        \item \key{SMP (symmetric multiprocessing)}: every processor is equal
          and runs the same OS; the scheduler can place any task on any CPU.
        \item \key{Asymmetric}: one master processor assigns work to slaves ---
          simpler, but the master is a bottleneck.
      \end{itemize}
    \item \key{Multicore}: several independent cores on one chip, each with its
      own registers and usually its own L1/L2 cache, sharing a last-level cache
      and memory.
  \end{itemize}
\end{frame}

\begin{frame}{Clusters, NUMA, and Why the OS Cares}
  \begin{itemize}
    \item \key{Cluster}: two or more complete systems connected by a network,
      sharing storage, presenting as one service. Used for \key{high
      availability} (failover) and \key{high-performance computing} (parallel
      work).
    \item \key{NUMA (non-uniform memory access)}: in large multiprocessors,
      memory is attached to nodes; a core accesses local memory faster than a
      remote node's memory. The scheduler and allocator should keep a task near
      its data.
    \item \key{Why it matters}: the OS must decide \emph{where} to run a task
      and \emph{where} to place its memory --- decisions that do not exist on a
      single-processor machine.
  \end{itemize}
  \vspace{0.2cm}
  \begin{alertblock}{The recurring theme}
    Every layer of hardware organization --- cache, memory, bus, processor ---
    becomes a policy question for the OS: which copy, which CPU, which node,
    which order.
  \end{alertblock}
\end{frame}

\begin{frame}{Comparison: SMP, Multicore, and Clusters}
  \centering
  \small
  \setlength{\tabcolsep}{4pt}
  \begin{tabular}{@{}>{\raggedright\arraybackslash}p{2.2cm}
      >{\raggedright\arraybackslash}p{3.1cm}
      >{\raggedright\arraybackslash}p{3.3cm}
      >{\raggedright\arraybackslash}p{3.9cm}@{}}
    \toprule
    & \textbf{SMP} & \textbf{Multicore} & \textbf{Cluster} \\
    \midrule
    What is shared? & Bus, clock, memory & Last-level cache, memory & Network, storage \\
    Failure unit & One processor & One core & One whole machine \\
    Main benefit & Throughput, scaling & Throughput, lower latency & Availability, scale-out \\
    OS challenge & Scheduling, locking & Cache affinity, locking & Distributed coordination \\
    \bottomrule
  \end{tabular}
  \vspace{0.35cm}
  \begin{itemize}
    \item All three give \key{parallelism}; only clusters survive the failure of
      an entire machine.
    \item All three make \key{concurrency} real for the OS --- the subject of
      Unit II.
  \end{itemize}
\end{frame}

\begin{frame}{Socratic Check 5}
  \begin{block}{Think, then answer}
    A machine has four cores but a single memory bus. Two programs each read a
    1~GB array. What limits the speedup, and what would you change first?
  \end{block}
  \vspace{0.2cm}
  \begin{itemize}
    \item The \key{memory bus} is the shared bottleneck: more cores do not add
      memory bandwidth. The workload is \key{memory-bound}, not CPU-bound.
    \item First fixes: improve \key{locality} (block the loops so each cache
      line is reused), then add cache or bandwidth --- not more cores.
    \item This is why the OS and compiler care about caches and NUMA placement,
      not just the number of processors.
  \end{itemize}
\end{frame}

% =============================================================================
% ACT 6 --- TAKEAWAYS AND CHECKS
% =============================================================================
\section{Takeaways and Checks}

\begin{frame}{Takeaways}
  \begin{itemize}
    \item A computer system is a \key{CPU, memory, controllers, and a bus} ---
      independent devices that run concurrently.
    \item Hardware gets the CPU's attention with an \key{interrupt}
      (asynchronous); the CPU gets the kernel's attention with a \key{trap}
      (synchronous).
    \item The \key{interrupt vector} dispatches to a handler; handlers must be
      short, hence top-half/bottom-half designs.
    \item Storage is a \key{hierarchy}: registers, cache, RAM, SSD, disk, tape
      --- faster and costlier per byte as you move up, kept coherent by
      \key{caching} and locality.
    \item I/O moves data by \key{polling}, \key{interrupts}, or \key{DMA}; the
      choice trades CPU work against latency, overhead, and bus pressure.
    \item \key{SMP, multicore, and clusters} replicate processors or machines;
      the OS must schedule and place work across them.
  \end{itemize}
\end{frame}

\begin{frame}{The Thread in One Picture}
  \begin{enumerate}
    \item \key{Keyboard interrupt}: the controller signals, the ISR queues a
      character, the shell runs.
    \item \key{Storage hierarchy}: \texttt{cat} asks for file blocks; the
      buffer cache checks RAM before touching the SSD or disk.
    \item \key{DMA}: the disk controller moves the block into memory; the CPU is
      interrupted once when it finishes.
    \item \key{Architecture}: the scheduler may run \texttt{cat} on another core
      while the DMA transfer proceeds --- and must keep the caches coherent.
  \end{enumerate}
  \vspace{0.2cm}
  \begin{block}{One sentence}
    The OS is an event-driven manager of a hierarchy of asynchronous devices ---
    interrupts tell it \emph{when}, storage tells it \emph{where}, and DMA tells
    it \emph{how} to move data without doing the moving itself.
  \end{block}
\end{frame}

\begin{frame}{Next: OS Operations and System Protection}
  \begin{itemize}
    \item Today: the hardware the OS manages and the signals it responds to.
    \item Next lecture: how the OS \emph{stays in control} --- dual-mode
      operation, the timer, privileged instructions, and the protection
      boundary that makes every mechanism in this lecture safe to share.
  \end{itemize}
  \vspace{0.3cm}
  \begin{block}{Reading}
    Silberschatz et al., \emph{Operating System Concepts}, Ch.~1 \S1.5
    (pp.~20--23) and Ch.~1 \S1.2 (pp.~6--18)
    \cite{SilberschatzOSC_operating_system_concepts}; Stallings,
    \emph{Operating Systems: Internals and Design Principles}, Ch.~1
    \S1.4--1.7 (pp.~15--38) \cite{StallingsOS_internals_design_principles}.
    Optional concrete depth: Bovet \& Cesati, \emph{Understanding the Linux
    Kernel}, Ch.~4 (pp.~131--188)
    \cite{BovetCesati_understanding_linux_kernel}.
  \end{block}
\end{frame}

\begin{frame}{Further Reading (Page-Cited)}
  \footnotesize
  \begin{itemize}
    \item A.~Silberschatz, P.~B.~Galvin, G.~Gagne, \emph{Operating System
      Concepts}, 8e --- Ch.~1 \S1.2 (computer-system organization, interrupts,
      storage, I/O, DMA, pp.~6--18), \S1.3 (architecture, pp.~15--20),
      \S1.7--1.8 (memory and storage management, pp.~24--26).
    \item W.~Stallings, \emph{Operating Systems: Internals and Design
      Principles}, 6e --- Ch.~1 \S1.4 (interrupts, p.~15), \S1.5 (memory
      hierarchy, p.~26), \S1.6 (cache, p.~29), \S1.7 (I/O techniques and DMA,
      p.~33).
    \item D.~P.~Bovet \& M.~Cesati, \emph{Understanding the Linux Kernel},
      3e --- Ch.~4 (interrupts and exceptions, pp.~131--188), optional.
  \end{itemize}
  \vspace{0.2cm}
  \muted{Page-verified via each book's \texttt{index.md} under
  \texttt{master\_supporting\_docs/CS502/supporting\_books/}.}
\end{frame}

\begin{frame}[allowframebreaks]{References}
  \footnotesize
  \bibliographystyle{plain}
  \bibliography{../../Bibliography_base}
\end{frame}

\end{document}
